\section{Localized crossings}\label{df:sec:crossings}

The main estimate in this section is Lemma~\ref{df:cr:rare-fast}: a crossing
confined to a strip of width of order $s$ is unlikely to have defect below
$-\varepsilon s$ when its length is much larger than $s$. The proof amplifies
any positive probability of such a crossing into an infinite concatenation
whose defect decreases linearly, contradicting the shape theorem.

We prepare three ingredients for this construction. The shape theorem first
gives short paths confined near their joining segments. A finite net then
makes this assertion simultaneous over all endpoints in two intervals.
Finally, an oriented-percolation argument joins sufficiently likely local
tests even when nearby tests overlap. Throughout the section, distances are
Euclidean unless indicated otherwise. A path constrained to a region has
all its vertices there, so its existence and cost are tested using only
edges with both endpoints in that region.

\begin{lemma}[Localized travel]\label{df:cr:local-travel}
Fix $D,\eta,\delta>0$. Uniformly over deterministic vertices $p,q\in\Z^2$
with $|q-p|\le Ds$, the probability tends to one as $s\to\infty$
that there is a path from $p$ to $q$ with cost at most
\[
  \mu(q-p)+\eta s
\]
whose vertices are within distance $\delta s$ of the segment $[p,q]$.
\end{lemma}

\begin{proof}
\emph{Uniform shape estimates.}
Use the constants $\alpha,\beta$ from \eqref{df:eq:norm-bounds}.
Equation~\eqref{df:eq:uniform-shape} implies, for each fixed $C>0$,
\begin{equation}\label{df:eq:cr-uniform-shape}
 \frac1s\max_{\substack{z\in\Z^2\\|z|\le Cs}}
       \abs{T(0,z)-\mu(z)}\longrightarrow0
 \qquad\text{almost surely}.
\end{equation}
Indeed, outside any sufficiently large fixed ball the error is bounded by an
arbitrarily small multiple of $|z|$, while the finitely many passage
times inside that ball contribute $o(s)$. The same shape theorem gives
\begin{equation}\label{df:eq:cr-escape}
 \PP\left(
   \inf_{\substack{z\in\Z^2\\|z|\ge\delta s/3}}
       T(0,z)\ge\alpha\delta s/6
       \right)\longrightarrow1.
\end{equation}
Translation invariance gives the same probability bounds at every
deterministic lattice origin. This uniformity will allow the origins to
depend on the endpoints and on $s$.

\emph{Short pieces cannot escape.}
Choose a fixed integer $B$ so large that $\beta D/B<\alpha\delta/24$.
Round the $B+1$ equally spaced points of $[p,q]$ to vertices
$p_0=p,p_1,\ldots,p_B=q$. Uniformly in the endpoints,
\[
 \mu(p_{j+1}-p_j)=\frac{\mu(q-p)}{B}+O(1).
\]
Apply \eqref{df:eq:cr-uniform-shape} and \eqref{df:eq:cr-escape} at these
finitely many deterministic vertices. With probability tending to one,
uniformly in $p,q$, every short passage time is at most
\[
 \frac{\mu(q-p)}B+
       \min\left\{\frac{\eta}{2B},\frac{\alpha\delta}{24}\right\}s,
\]
and escape from the ball of radius $\delta s/3$ about its starting vertex
costs at least $\alpha\delta s/6$. For each $j=0,\ldots,B-1$, choose a path
from $p_j$ to $p_{j+1}$ with cost at most $T(p_j,p_{j+1})+1$.
For all sufficiently large $s$, each chosen path has cost strictly below
the corresponding escape bound. Nonnegativity of the weights forces it
to stay inside the ball about $p_j$: its initial portion up to a vertex
outside would already cost at least the escape bound.

\emph{Concatenation.}
Concatenating the $B$ paths gives total cost
at most $\mu(q-p)+\eta s/2+B\le\mu(q-p)+\eta s$. Each ball is centered
within bounded distance of the original segment, so it lies in its
$\delta s$-neighborhood for large $s$. All probability bounds used above
are uniform in the deterministic vertices, proving the stated uniformity.
\end{proof}

The next form permits the endpoints to be selected after the environment
is observed. The intervals themselves remain deterministic, and the
assertion holds for every endpoint pair on one event of high probability.

\begin{lemma}[Travel between intervals]\label{df:cr:interval-travel}
Fix $D,\eta,\delta>0$. Let $I_1,I_2$ be deterministic vertical intervals
on lattice columns, and suppose that the diameter of $I_1\cup I_2$ is at
most $Ds$. Uniformly over their locations, with probability tending to
one as $s\to\infty$, every pair
$p\in I_1\cap\Z^2$, $q\in I_2\cap\Z^2$ can be joined by a path of
cost at most $\mu(q-p)+\eta s$ within distance $\delta s$ of $[p,q]$.
\end{lemma}

\begin{proof}
The assertion is vacuous if either interval contains no vertex.

\emph{Finite nets and vertical attachments.}
Let $\beta$ be the upper norm bound in \eqref{df:eq:norm-bounds}, and
choose a fixed $\gamma>0$ so small that
\[
 \left(\frac{4\kappa}{\lambda}+2\beta\right)\gamma\le\eta/2,
 \qquad \gamma\le\delta/2.
\]
In each interval, choose a net of lattice vertices containing its extreme
lattice vertices and having consecutive gaps at most
$\lfloor\gamma s\rfloor$. For all large $s$, each net needs at most
$2D/\gamma+2$ vertices. Assign every interval vertex to a net vertex at
vertical distance at most $\gamma s$, using only edges of its intervening gap.

For any such gap, the total weight $V$ is a sum of at most
$\lfloor\gamma s\rfloor$ independent edge weights.
Lemma~\ref{df:ga:moment}, with $A=\gamma$ and $v=s$, gives
\[
 \PP\left(V>\frac{2\kappa}{\lambda}\gamma s\right)
 \le e^{-\kappa(1-\log2)\gamma s}.
\]
There are only boundedly many gaps. Thus, with probability tending to
one, all vertical connections from interval vertices to their chosen
net vertices cost at most $2\kappa\gamma s/\lambda$.

\emph{Connections between net vertices.}
Apply Lemma~\ref{df:cr:local-travel}, with errors $\eta/2$ and $\delta/2$,
to every pair of net vertices. Its uniformity and the bounded number
of pairs give these connections simultaneously.

\emph{All endpoint pairs on the same event.}
On the intersection of these high-probability events, take arbitrary
$p\in I_1\cap\Z^2$ and $q\in I_2\cap\Z^2$. Concatenate their vertical
attachments with the path between the assigned net vertices
$\widehat p,\widehat q$. The resulting cost is at most
\[
 \begin{split}
 \mu(\widehat q-\widehat p)+\eta s/2
       +\frac{4\kappa}{\lambda}\gamma s
 &\le \mu(q-p)+\eta s/2
       +\left(\frac{4\kappa}{\lambda}+2\beta\right)\gamma s.
 \end{split}
\]
The segments $[p,q]$ and $[\widehat p,\widehat q]$ have Hausdorff distance
at most $\gamma s$. The choice of $\gamma$ therefore gives both the cost
bound and the required localization, including the vertical attachments.
The same event works for every $p,q$, and its probability bound is uniform
in the deterministic locations of the intervals.
\end{proof}

We will place local crossing and connection tests on an auxiliary copy of
$\Z^2$. The following elementary lemma needs only a uniform bound on their
failure probabilities and independence at sufficiently separated sites.

\begin{lemma}[High-density oriented paths]\label{df:cr:oriented-paths}
For every integer $k_0\ge0$ there is $\zeta=\zeta(k_0)>0$ with the
following property. Let $(G_v)_{v\in\Z^2}$ be events such that any finite
family indexed by sites at pairwise sup distance greater than $k_0$ is
independent. If $\PP(G_v^c)\le\zeta$ for every $v$, then with positive
probability there is an infinite path starting at $0$, using steps
$(1,0)$ and $(0,1)$, all of whose sites satisfy $G_v$.
\end{lemma}

\begin{proof}
\emph{A finite reachable set if there is no infinite path.}
Call a site good when $G_v$ occurs. Let $S$ consist of the endpoints of
finite oriented paths from $0$ for which every site except possibly the
last is good. In particular $0\in S$. If $S$ were infinite, the finitely
branching tree of these paths would have an infinite branch; every site
on that branch would be good. Hence, on the event of no infinite good
path, $S$ is finite.

\emph{A surrounding contour forces bad sites.}
For such a finite $S$, orient the boundaries of the unit squares centered
at its sites counterclockwise and cancel shared sides. The remaining
edges form a balanced directed graph and hence decompose into closed
directed trails, with no edge repeated in a trail. Repeated vertices are
allowed, so corner contacts cause no difficulty. Every edge has an
$S$-square on its left and a complementary square on its right.
Winding number is additive under cancellation and decomposition. Its
total about $0$ is one, contributed by the square centered at $0$;
therefore some trail has nonzero winding.

Let $m$ be the length of this trail. Closedness implies that exactly
$m/2$ of its steps are upward or leftward. For these steps, the site on
the right is respectively an east or north neighbor of the site on the
left. The left-hand site must be bad: if it were good, extending a path
that reaches it would put this neighbor in $S$ as well.

\emph{Independent bad sites along a fixed contour.}
A site borders at most four edges of the trail, so at least $m/8$
distinct sites are forced to be bad. For any fixed candidate trail these
sites are deterministic. Greedy selection in a fixed order gives at least
\[
 \frac{m}{8(2k_0+1)^2}
\]
of them at pairwise sup distance greater than $k_0$: each selected site
removes at most $(2k_0+1)^2$ candidates. Their tests are independent, so
the probability of all the badness conditions forced by this particular
trail is at most $\zeta^{m/[8(2k_0+1)^2]}$.

\emph{Summing over contours.}
Nonzero winding forces $0$ into the coordinate bounding box of the
trail. A length-$m$ trail has coordinate spans at most $m$, so it lies
in a square of side $O(m)$ about $0$. Choosing its starting vertex and
then its step sequence bounds the number of candidates by
$C(m+1)^2 4^m$. Consequently the probability of no infinite good path
is at most
\[
 \sum_{m\ge4} C(m+1)^2
       \left(4\zeta^{1/[8(2k_0+1)^2]}\right)^m.
\]
This is less than one when $\zeta$ is sufficiently small. The argument
uses no stationarity assumption on the tests.
\end{proof}

We now combine these ingredients. The relevant defect is measured against
one of the extremal supports $f_\sigma$. Its one-sided contact with $\mu$
allows successive crossings to be connected with defect as small a multiple
of $s$ as needed, by drifting consistently to the corresponding side of
the direction $u$.

\begin{lemma}[Rare fast localized crossings]\label{df:cr:rare-fast}
Fix $A,R,\varepsilon>0$ and $\sigma\in\{-,+\}$. As $s\to\infty$ and
$n/s\to\infty$, with $s$ real and $n$ a positive integer,
\begin{equation}\label{df:eq:cr-rare-fast}
 \PP\left(
 \begin{array}{l}
 \text{there is a path $P$ from $x=0$ to $x=n$ with}
       \ H_\sigma(P)<-\varepsilon s,\\
 \text{all of whose vertices satisfy }
       -An\le x\le(1+A)n,\quad |r|\le Rs
 \end{array}\right)\longrightarrow0.
\end{equation}
The same assertion holds, with the same probability, for every exact
lattice translate of the event.
\end{lemma}

\begin{proof}
Suppose that the probability in \eqref{df:eq:cr-rare-fast} is at least
$p_0>0$ along a sequence satisfying the stated limits. At each auxiliary
site we will try several disjoint translates of this crossing event and
require connections to both forward neighbors. Enough trials make a
successful crossing likely; the one-sided contact estimate makes the
connections likely without spending the whole crossing deficit.

\emph{1. Parameters and site geometry.}
The dependence radius must be fixed before we choose how many trials to
make. Set
\[
 k_0=\left\lceil4(A+3)+8\right\rceil+1
\]
and let $\zeta$ be supplied by Lemma~\ref{df:cr:oriented-paths}. Choose a
fixed integer $M$ so large that $(1-p_0)^M<\zeta/2$. For each member of
the sequence put
\[
 S_0=100M(R+1)\lceil s\rceil,
 \qquad d=\lceil D_0s\rceil,
\]
where the fixed constant $D_0$ will be chosen below, after $M$.

Identify the signs $+$ and $-$ with $1$ and $-1$ when they multiply
coordinates. For a test site $v=(v_1,v_2)\in\Z^2$, write
\[
 i=v_1+v_2,
 \qquad X_i=i(n+d),
 \qquad Y_v=\sigma(i+v_2)S_0.
\]
An oriented step sends $i$ to $i+1$ and changes $Y_v$ by $\sigma S_0$
or $2\sigma S_0$. Thus every step advances to the next longitudinal
block and moves transversely to the same prescribed side.

\emph{2. Independent crossing trials within a site.}
At site $v$, make $M$ trials using exact lattice translates of the event
in \eqref{df:eq:cr-rare-fast}, all with initial column $X_i$, and with their
strips disjoint and contained in $|r-Y_v|\le S_0/4$. Here is an explicit
choice, which also accounts for lattice phases. Put
\[
 h_j=(2j-M-1)\,3(R+1)\lceil s\rceil,
 \qquad
 t_{v,j}=\left(X_i,
    \left\lfloor bX_i+Y_v+h_j+\tfrac12\right\rfloor\right),
 \quad 1\le j\le M.
\]
The translated strip has center $r(t_{v,j})$, within $1/2$ of
$Y_v+h_j$. For $s\ge1$, consecutive centers are separated by more than
$2Rs$, and every strip of half-width $Rs$ lies inside the designated
$S_0/4$ window. Thus the trials use disjoint internal edge sets.
Translation leaves both displacement and its supporting functional
unchanged, so each trial has probability at least $p_0$, exactly as in
the original event. Their probability of all failing is at most
$(1-p_0)^M$.

\emph{3. Connection tests for every endpoint pair.}
For each of the two forward neighbors $v'$ of $v$, impose an additional
test. Every vertex on
\[
 x=X_i+n,\quad |r-Y_v|\le S_0/4
\]
must be connectible to every vertex on
\[
 x=X_{i+1},\quad |r-Y_{v'}|\le S_0/4
\]
by a path of cost at most $f_\sigma$ of its displacement plus
$\varepsilon s/2$, lying in
\begin{equation}\label{df:eq:cr-connector-region}
 X_i+n-s\le x\le X_{i+1}+s,
 \qquad |r-Y_v|\le4S_0.
\end{equation}
These are existence tests using only the internal edges of this region.
Requiring every endpoint pair will permit the successful trials to be
chosen later, after the environment has been observed.

\emph{4. Choosing the connector length.}
We next choose $D_0$ and verify that both connection tests pass with
probability tending to one uniformly in $v$. A connector displacement
$z$ satisfies
\[
 x(z)=d,
 \qquad S_0/2\le\sigma r(z)\le5S_0/2\le3S_0.
\]
Set $C_0=200M(R+1)$, so that $S_0/s\le C_0$ for $s\ge1$.
The one-sided expansion \eqref{df:eq:support-contact} supplies $h_0>0$
such that
\[
 0\le\mu(u+h e_2)-f_\sigma(u+h e_2)
       \le\frac{\varepsilon}{12C_0}|h|
 \quad\text{if }0<\sigma h\le h_0.
\]
Choose $D_0\ge1$ with $3C_0/D_0\le h_0$. For every possible connector,
$h=r(z)/d$ is on the required side and has absolute value at most $h_0$.
Homogeneity therefore gives
\[
 \mu(z)-f_\sigma(z)
   \le\frac{\varepsilon}{12C_0}|r(z)|\le\varepsilon s/4.
\]
This leaves another $\varepsilon s/4$ for the travel error.

The two endpoint intervals have union diameter at most $Ds$ for a fixed
$D$ depending only on $M,R,D_0$. Their joining segments have longitudinal
coordinates between the two endpoint columns and satisfy
$|r-Y_v|\le9S_0/4$. Since $|b|\le1$, moving Euclidean distance
$\delta s$ changes $x$ by at most $\delta s$ and $r$ by at most
$\sqrt2\delta s$. Thus a tube with any fixed $0<\delta\le1/2$ fits
inside \eqref{df:eq:cr-connector-region}. Lemma~\ref{df:cr:interval-travel},
with cost error $\varepsilon s/4$, proves the required connection assertions
simultaneously for all endpoint pairs. Its probability estimate is
uniform in the intervals' locations, including their lattice phases.

\emph{5. A dependence radius independent of the later choices.}
Call $v$ good when at least one of its trials succeeds and both connection
tests pass. Its bad probability is at most
$(1-p_0)^M+o(1)$, uniformly in $v$. We verify that the initially chosen
$k_0$ is a dependence radius. Take a sufficiently large member of the
sequence that $n\ge d+s$. The entire test at $v$ is supported on internal
edges in
\[
 |x-X_i|\le(A+3)n,
 \qquad |r-Y_v|\le4S_0.
\]
To compare two sites, use the integer coordinates $i=v_1+v_2$ and
$j=i+v_2$. If their support regions share a vertex, the longitudinal
and transverse bounds give, respectively,
\[
 |\Delta i|\le2(A+3),
 \qquad |\Delta j|\le8,
\]
where $\Delta$ denotes the difference between the two sites.
Since $v_2=j-i$ and $v_1=2i-j$, this implies
\[
 |\Delta v_2|\le2(A+3)+8,
 \qquad |\Delta v_1|\le4(A+3)+8<k_0.
\]
Tests at sites at pairwise sup distance greater than $k_0$ therefore use
pairwise disjoint edge sets and are jointly independent. In particular,
the dependence radius does not increase when $M$ or $D_0$ is chosen.
The order of choices was consequently legitimate: $k_0$ and $\zeta$
first, then $M$, then $D_0$, and finally a sufficiently large member of
the sequence. For that member the bad probability is at most $\zeta$.

\emph{6. Concatenation at one fixed scale.}
By Lemma~\ref{df:cr:oriented-paths}, with positive probability there is an
infinite oriented path of good sites starting at $0$. Fix this one scale
$s,n,d,S_0$ and a configuration with such a path. At each visited site
choose a successful trial and join it to the next by a tested connector.
Prepend a vertical path from the lattice origin to the initial trial's
starting vertex; its defect $H_\sigma$ is a finite random quantity $C$.
If $z_t$ is the endpoint of the trial at the $t$th site, with the initial
site numbered $0$, additivity gives a path from $0$ to $z_t$ of defect
at most
\begin{equation}\label{df:eq:cr-linear-deficit}
 C-(t+1)\varepsilon s+t\varepsilon s/2.
\end{equation}
The trials and connectors need not be disjoint: costs are counted with
edge multiplicities, so their concatenation still bounds passage time
from above.

\emph{7. The shape-theorem contradiction.}
At that site, $i=t$ and $0\le v_2\le t$. Hence
\[
 x(z_t)=t(n+d)+n,
 \qquad |r(z_t)|\le2tS_0+S_0/4.
\]
Thus $|z_t|\to\infty$ and $|z_t|\le C'(t+1)$ for a deterministic
constant $C'$ at this fixed scale. The limit here is $t\to\infty$;
the parameters $s,n,d,S_0$ remain fixed. On the probability-one event
of the origin shape theorem, its uniform outside-ball estimate applies
also to these adaptively chosen endpoints and yields
\[
 T(0,z_t)\ge\mu(z_t)-o(t)\ge f_\sigma(z_t)-o(t),
\]
where the supporting inequality is from Lemma~\ref{df:prelim:supports}.
This contradicts \eqref{df:eq:cr-linear-deficit}, whose negative part grows
at least as $t\varepsilon s/2$. The infinite-path event consequently has
probability zero, contradicting its construction. No sequence with
crossing probability bounded below by $p_0$ exists, proving the asserted
joint limit. Exact lattice-translation invariance proves the last assertion.
\end{proof}
